% !PRINTED 374
\weberpage{374}
6. \gesperrt{The order $O$ belonging to the conductor $Q$ is called composed or compounded from $O_1$ and $O_2$.}

We set symbolically
\weq{2}{O = O_1 O_2.}
The discriminants of these three orders are
\weq{3}{D = Q^2 \Delta, \quad D_1 = Q_1^2 \Delta, \quad D_2 = Q_2^2 \Delta.}
If $Q_1$ and $Q_2$ are relatively prime, then the order composed from the two is the maximal order $O_0$.

To carry out the composition of the orders one can proceed just as in the preceding section. We set
\weq{4}{Q_1 = Qm_1, \quad Q_2 = Qm_2,}
and then $m_1$, $m_2$ are relatively prime to one another.

Let $\mathfrak{a}$ and $\mathfrak{b}$ then denote two ideals which are relatively prime to $Q_1 Q_2$.

If then
\[
\begin{aligned}
\mathfrak{a}&\ \text{equivalent}\ \mathfrak{a}'&&\ \text{with respect to}\ O_1,\\
\mathfrak{b}&\ \text{\glqq}\ \mathfrak{b}'&&\ \text{\glqq}\ O_2,
\end{aligned}
\]
then
\[
\begin{aligned}
\frac{\mathfrak{a}}{\mathfrak{a}'}&=\eta_1\ \text{a number in}\ O_1,\\
\frac{\mathfrak{b}}{\mathfrak{b}'}&=\eta_2\ \text{\glqq}\ \text{\glqq}\ \text{\glqq}\ O_2,
\end{aligned}
\]
i.\,e. $\eta_1$ is congruent to a rational number with respect to the modulus $Q_1$, and $\eta_2$ with respect to the modulus $Q_2$. Hence both, and consequently their product, are congruent to a rational number with respect to the modulus $Q$ and belong to the order $Q$. % [sic] German source prints Q instead of O.
Accordingly we have also
\[
\mathfrak{c} = \mathfrak{a}\mathfrak{b} \ \text{equivalent to} \ \mathfrak{c}' = \mathfrak{a}'\mathfrak{b}' \ \text{with respect to} \ O.
\]
Let $a_1$ be the least natural number divisible by $\mathfrak{a}$, and let $\mathfrak{b}$ be relatively prime to $a_1$ (and consequently also to $\mathfrak{a}$); and if $a_2$ is the least natural number divisible by $b$, then $a_1$ and $a_2$ are also relatively prime. % [sic] The German source prints the second b in roman type; context indicates the ideal $\mathfrak{b}$.
One can then always, whatever $b_1$, $b_2$ may be, satisfy the congruences
\[
\begin{aligned}
b_1 &\equiv m_1 b \ (\mathrm{mod}\ 2a_1),\\
b_2 &\equiv m_2 b \ (\mathrm{mod}\ 2a_2)
\end{aligned}
\]
and accordingly assume the basis forms of $\mathfrak{a}$ and $\mathfrak{b}$ in the shape:
%CONT>>
